% =====================================================================
% Capítulo 4 — Reconstrucción de PET de dosis baja
% =====================================================================
\chapter{Modelos de difusión para reconstrucción de escaneos PET}
\label{cap:pet}

Este capítulo desarrolla la contribución principal del trabajo: un
sistema completo de reconstrucción de PET de dosis baja. El problema
se formuló en la \Cref{subsec:pet-problema-inverso} como un problema
inverso —recuperar el escaneo de dosis completa $\mathbf{x}$ a partir
de la observación degradada $\mathbf{y}$, adquirida con un veinteavo
de la dosis— y la \Cref{subsec:difusion-condicional} presentó las dos
familias de soluciones basadas en difusión que coexisten en la
literatura: el condicionamiento supervisado y el prior no condicional
con guiado de medición. Siguiendo la trayectoria que dibuja el estado
del arte (\Cref{sec:estado-del-arte}), se decide construir
\textbf{ambos enfoques en paralelo} y compararlos sobre el mismo
conjunto de test, añadiendo además dos líneas base discriminativas
que delimitan la aportación real del paradigma de difusión. El diseño
persigue dos objetivos: ser \emph{defendible científicamente} —cada
decisión está justificada por la literatura o por los fundamentos
teóricos del \Cref{cap:marco-teorico}—.

El capítulo sigue el ciclo de vida natural del sistema: análisis del
conjunto de datos (\Cref{sec:pet-datos}), preprocesamiento y formato
final (\Cref{sec:pet-preprocesado}), metodología de entrenamiento
—incluido el entorno de ejecución en la nube—
(\Cref{sec:pet-entrenamiento}), arquitectura y enfoque de cada
\textit{pipeline} (\Cref{sec:pet-arquitectura}), métricas de
evaluación (\Cref{sec:pet-metricas}), líneas base de referencia
(\Cref{sec:pet-baselines}), comparación de resultados
(\Cref{sec:pet-resultados}) y análisis de costes
(\Cref{sec:pet-costes}).

% ---------------------------------------------------------------------
\section{El conjunto de datos: análisis de los volúmenes NIfTI}
\label{sec:pet-datos}

Los datos proceden del \textit{Ultra-low Dose PET Imaging
Challenge} (UDPET)~\cite{udpet2025}, un desafío internacional que
publica pares de escaneos PET de cuerpo entero adquiridos con dosis
completa y reconstruidos también a partir de fracciones reducidas de
los eventos detectados. El conjunto empleado consta de \textbf{371
pares} de volúmenes: para cada paciente se dispone de la versión de
dosis completa y la versión con dosis reducida a un veinteavo (1/20).

Los volúmenes se distribuyen en formato \textbf{NIfTI}
(\textit{Neuroimaging Informatics Technology Initiative}), el formato
estándar de intercambio en neuroimagen e imagen médica volumétrica:
cada fichero \texttt{.nii.gz} empaqueta el \textit{array}
tridimensional de intensidades junto con una matriz afín de
$4 \times 4$ que sitúa los vóxeles en el espacio físico del escáner.
El análisis exploratorio de los volúmenes arroja las siguientes
características:

\begin{itemize}[leftmargin=*]
    \item Cada volumen tiene forma $(440, 440, 644)$ —cortes axiales
    de $440 \times 440$ apilados a lo largo del eje craneocaudal—
    con vóxeles isotrópicos de $1{,}65$~mm.
    \item Se asume que cada fichero corresponde a un paciente
    independiente.
\end{itemize}

Las estadísticas de intensidad relevantes para el diseño del
preprocesamiento, calculadas sobre los volúmenes del dataset son:

\begin{itemize}[leftmargin=*]
    \item La distribución de intensidades en \textit{foreground}
    (vóxeles con señal) es muy asimétrica: mediana de 251 cuentas,
    percentil 99 en torno a 9.600, percentil 99,9 en torno a 19.000
    y máximo cercano a $1{,}3 \times 10^{5}$. Es la cola pesada
    típica de una señal de recuento con regiones calientes (vejiga,
    corazón, lesiones).
    \item Aproximadamente el 18\,\% de los vóxeles del volumen son
    \textit{foreground} (valor estrictamente positivo); el resto es
    fondo exactamente nulo, fuera del cuerpo del paciente.
    \item La desviación típica del residuo \texttt{low\,-\,full}
    restringido al \textit{foreground} es del orden del 31\,\% de la
    media de señal, lo cual cuantifica el nivel de degradación
    introducido por la reducción de dosis —el ruido Poisson de
    recuento descrito en la \Cref{sec:pet-ruido}—.
\end{itemize}

% ---------------------------------------------------------------------
\section{Preprocesado y formato final}
\label{sec:pet-preprocesado}

El preprocesamiento convierte cada par de volúmenes NIfTI en una
colección de cortes axiales 2D normalizados, listos para el
entrenamiento. Las decisiones, en orden de aplicación, son las
siguientes.

\paragraph{Unidad espacial.}
Se trabaja con \textit{slices} axiales en 2D puro. Un U-Net 3D
sobre volúmenes de $440^2 \times 644$ es inviable con el hardware
disponible; las variantes 2.5D (canales adicionales con
\textit{slices} contiguos) se reservan como extensión natural si la
formulación 2D muestra inconsistencias inter-\textit{slice}
(\Cref{sec:lineas-futuras}).

\paragraph{Recorte y reescalado.}
Para cada volumen se calcula una \textit{bounding box} única a
partir del \textit{foreground} de la versión de \textbf{dosis
baja}, y se aplica idéntica al volumen pareado de dosis completa,
garantizando que ambos comparten exactamente el mismo recorte.
La escala de normalización $M$ y el filtrado de \textit{slices} se calculan todos a partir de la observación de
dosis baja, la única señal disponible en inferencia para un volumen no visto. Cada \textit{slice} axial recortado se
redimensiona a $256 \times 256$ píxeles mediante interpolación bilineal.

La resolución de trabajo $256^2$ es menor que la nativa
($440 \times 440$). No es posible entrenar directamente sobre
$440^2$: la arquitectura U-Net empleada tiene seis bloques —cinco
operaciones de \textit{down}/\textit{upsampling}, $2^5 = 32$—, de
modo que el lado de la imagen de entrada debe ser divisible entre
32, y $440$ no lo es. Tamaños como $384$ o $416$ —más cercanos a
la calidad nativa y divisibles entre 32— sí serían admisibles,
pero no se han explorado: el coste de cualquier fase del
entrenamiento crece de forma aproximadamente cuadrática con el
lado de la imagen, y ya resulta elevado a $256^2$. La elección de
$256^2$ equilibra fidelidad anatómica y coste computacional dentro
del alcance del proyecto.

\paragraph{Normalización por-volumen estabilizadora de varianza.}
Las cuentas se mapean al rango $[-1, 1]$ mediante la
transformación
\begin{equation}
    \mathbf{x}' = 2 \cdot
        \frac{\operatorname{arcsinh}(\mathbf{x}/k)}
             {\operatorname{arcsinh}(M/k)}
        - 1,
    \label{eq:asinh-norm}
\end{equation}
donde $k = 10$ actúa como codo suave entre el régimen lineal y
el logarítmico, y $M$ es el percentil 99,5 del \textbf{volumen de
dosis baja}. El mismo $M$ derivado de la dosis baja se aplica
también al volumen pareado de dosis completa, de modo que ambos
comparten escala y siguen siendo directamente comparables tras la
normalización. Calcular $M$ exclusivamente sobre la dosis baja es
esencial para no introducir fuga del objetivo: $M$ debe poder
recomputarse en inferencia a partir de la única señal disponible,
el escaneo de dosis baja de un volumen no visto. La función
$\operatorname{arcsinh}$ es la transformación canónica
estabilizadora de varianza para señales positivas de cola pesada
con estructura aproximadamente Poisson: se comporta linealmente
cerca de cero y logarítmicamente para valores grandes, sin la
patología de $\log(0)$. Los parámetros $(M, k)$ se almacenan por
volumen para invertir la transformación durante la evaluación.

\begin{figure}[h!]
    \centering
    \includegraphics[width=0.92\linewidth]{img/asinh_normalization.pdf}
    \caption{Normalización per-volumen estabilizadora de varianza
    (\Cref{eq:asinh-norm}) con los parámetros reales de un volumen del
    conjunto ($M = 11\,827$ cuentas, percentil 99,5; $k = 10$). La
    curva central traza el mapeo $x \mapsto x'$: el origen va a $-1$ y
    el percentil de referencia $x=M$ a $+1$; el recuadro ampliado
    muestra el régimen casi lineal cerca del origen y el codo en $x=k$,
    donde la transformación pasa al régimen logarítmico. El histograma
    superior es la distribución de cuentas del \textit{foreground}
    —asimétrica y de cola pesada, concentrada cerca de cero—; el
    lateral es esa misma distribución tras la transformación, ya
    repartida de forma estable sobre la banda $[-1, +1]$ en la que
    opera la difusión. El $0,5\%$ de vóxeles más brillantes sobresale
    ligeramente por encima de $+1$ sin \textit{clipping}, lo que
    mantiene exacta la inversión.}
    \label{fig:asinh-norm}
\end{figure}

\paragraph{Elección del percentil de referencia.}
La elección de $M$ como percentil 99,5 —y no como el máximo
(percentil 100)— es deliberada. Si se fijara $M$ al valor más alto
del volumen, un único vóxel atípico muy brillante podría
desviarse mucho de la distribución y aplanar el resto de la
imagen, perdiendo detalle en la masa de tejido. Los vóxeles que quedan bastante
por encima de $+1$ están efectivamente fuera de distribución para
el proceso de ruido hacia delante, y el modelo tiende a
reconstruirlos peor (saturación suave / subestimación de la
intensidad máxima). No se \textit{clippean} esos valores —dejarlos
sobresalir por encima de $+1$ mantiene exacta la inversión de la
transformación—, pero su proporción debe ser pequeña. Un análisis
del porcentaje de \textit{foreground} mapeado por encima de $+1$
para distintos percentiles arroja:
\begin{center}
\begin{tabular}{lc}
\toprule
Percentil & \% \textit{foreground} mapeado por encima de $+1$ \\
\midrule
90{,}0 & 10{,}00\,\% \\
95{,}0 & 5{,}00\,\% \\
99{,}0 & 1{,}00\,\% \\
\textbf{99{,}5} & \textbf{0{,}50\,\%} (por defecto) \\
99{,}9 & 0{,}10\,\% \\
\bottomrule
\end{tabular}
\end{center}
Bajar el percentil dedicaría más del rango $[-1, +1]$ a la masa
de vóxeles —lo cual podría mejorar las métricas estructurales en
tejido blando (SSIM, contraste de intensidad media)— a costa de
empeorar la fidelidad de intensidad máxima en regiones calientes
y las métricas de preservación de intensidad en espacio de
cuentas. El percentil 99,5 deja fuera únicamente un 0,5\,\% de
\textit{foreground} —\textit{outliers} muy brillantes—, lo cual
preserva la distribución original de los datos a la vez que los
normaliza de forma efectiva.

\paragraph{Filtrado de \textit{slices}.}
Tras el recorte y la normalización, se conservan únicamente
aquellos \textit{slices} axiales cuyo \textit{foreground} supera
un umbral mínimo de ocupación ($1\%$ de píxeles por encima de un
valor bajo). Los extremos vacíos del campo de vista se descartan para no
penalizar al modelo con \textit{slices} triviales.

\paragraph{Ausencia de aumento de datos.}
El PET presenta asimetría clínica relevante —corazón en el lado
izquierdo, hígado en el derecho, etcétera—; un volteo horizontal
aleatorio induciría una lateralidad incorrecta en el modelo
aprendido. El tamaño efectivo del conjunto de entrenamiento hace
además innecesaria la augmentación para evitar sobreajuste.

\paragraph{Particiones y formato final.}
Los 371 volúmenes se dividen \emph{a nivel de paciente} en
proporción $80 / 10 / 10$ —296 volúmenes de entrenamiento, 37 de
validación y 38 de test— con semilla fija, y la asignación se
persiste en un fichero JSON compartido por todos los
\textit{pipelines} y líneas base, de modo que todos los modelos ven
exactamente las mismas particiones. El preprocesamiento se ejecuta
una única vez en modo \textit{offline}; cada \textit{slice}
procesado se almacena como tensor PyTorch \texttt{.pt} en precisión
\texttt{float16} para evitar la decodificación repetida de NIfTI en
cada época. El resultado final es una caché de \textbf{477.690
ficheros} \texttt{.pt} —238.845 pares de cortes
dosis baja / dosis completa, unos 644 cortes conservados por
paciente— con un tamaño total aproximado de 60~GB. De ellos,
aproximadamente $1{,}9 \times 10^{5}$ pares pertenecen al conjunto
de entrenamiento.

% ---------------------------------------------------------------------
\section{Metodología de entrenamiento}
\label{sec:pet-entrenamiento}

Esta sección reúne \emph{toda} la operativa de entrenamiento del
proyecto: la configuración numérica y las técnicas de estabilización
y aceleración (\Cref{subsec:pet-config}), y el entorno de ejecución
profesional construido sobre la nube para ejecutar los experimentos
de forma reproducible (\Cref{sec:pet-entorno-cloud}). Las secciones
posteriores se remiten aquí cuando necesitan algún detalle del
entrenamiento.

\subsection{Configuración del entrenamiento}
\label{subsec:pet-config}

El paradigma de difusión combina tres elecciones que mejoran sobre
la línea base del experimento de mariposas
(\Cref{sec:mariposas-paradigma}), todas ellas definidas formalmente
en el marco teórico:

\begin{itemize}[leftmargin=*]
    \item \textbf{Predicción de velocidad}
    (\textit{v-prediction}, \Cref{eq:v-def}), numéricamente estable
    en los dos extremos del rango de \textit{timesteps}
    (\Cref{subsec:parametrizaciones}). Su activación en
    \texttt{diffusers} consiste en fijar
    \texttt{prediction\_type="v\_prediction"} en el
    \textit{scheduler}.
    \item \textbf{Calendario cosenoidal} de ruido
    (\Cref{eq:cosine-schedule}), con $T = 1000$ \textit{timesteps}
    de entrenamiento.
    \item \textbf{Pérdida MSE uniforme sobre $\mathbf{v}$}
    (\Cref{eq:loss-v}), con el canal de condicionamiento presente u
    omitido según el \textit{pipeline}
    (\Cref{sec:pet-arquitectura}). Se opta por ponderación uniforme
    sobre los \textit{timesteps}, dejando el reponderado
    \textit{Min-SNR-$\gamma$}~\cite{hang2023minsnr} como mejora
    candidata si la convergencia resultara insuficiente.
\end{itemize}

Los hiperparámetros de optimización son:

\begin{itemize}[leftmargin=*]
    \item Optimizador \texttt{AdamW} con tasa de aprendizaje
    $\eta = 1{,}4 \times 10^{-4}$ y \textit{cosine warmup} de
    500 pasos. El valor de $\eta$ está escalado por
    $\sim\sqrt{32/16}$ sobre la línea base $10^{-4}$ del modelo de
    mariposas, para acompañar el incremento de \textit{batch} sin
    alterar la dinámica del entrenamiento.
    \item \textit{Batch} de entrenamiento de 32 \textit{slices} con
    acumulación de gradiente unitaria (\textit{batch} efectivo de
    32). En el \textit{fallback} local sobre MPS se emplea
    \textit{micro-batch} de 4 con acumulación de 4 pasos.
    \item Recorte de gradientes a norma máxima $1{,}0$.
    \item Se entrena durante \textbf{100 épocas}, lo cual equivale a
    aproximadamente $6 \times 10^{5}$ pasos de gradiente sobre los
    $\sim 1{,}9 \times 10^{5}$ pares de cortes de entrenamiento.
\end{itemize}

Sobre esta base se aplican cuatro técnicas de estabilización y
aceleración:

\paragraph{Media móvil exponencial de los pesos (EMA).}
Se mantiene una copia sombra de los pesos del U-Net actualizada como
media móvil exponencial con factor de decaimiento $0{,}9999$. Los
pesos EMA —un promedio suavizado de la trayectoria de optimización—
son los que se evalúan y se exportan como \textit{checkpoint} de
inferencia; aportan típicamente $0{,}5$--$1$~dB adicionales de PSNR
sin coste de entrenamiento. El estado EMA se registra junto al del
optimizador, de modo que la reanudación de un entrenamiento
interrumpido lo recupera intacto.

\paragraph{Precisión mixta \texttt{bf16}.}
El cómputo se ejecuta en \texttt{bfloat16} sobre los \textit{tensor
cores} de la GPU, prácticamente el doble de \textit{throughput} que
\texttt{float32} sin riesgo apreciable de desbordamiento numérico
(el formato conserva el rango del \texttt{float32}). Se habilita
además el modo TF32 para las multiplicaciones residuales en
\texttt{float32}. En el \textit{fallback} local sobre MPS se
mantiene \texttt{float32} puro, por las inestabilidades conocidas
del \textit{backend} con \texttt{float16}.

\paragraph{Compilación del modelo con \texttt{torch.compile}.}
La pasada hacia delante del U-Net se compila con TorchInductor, el
compilador JIT de PyTorch, que fusiona kernels y elimina
lanzamientos redundantes de GPU; en un U-Net limitado por cómputo
como el empleado, la mejora observada es de entre $1{,}3\times$ y
$2\times$ en velocidad de entrenamiento. La compilación se aísla en
un manejador exclusivo de la pasada hacia delante que comparte
parámetros con el modelo original, de modo que el optimizador, la
EMA, el recorte de gradientes y el guardado/carga de
\textit{checkpoints} siguen operando sobre el modelo sin compilar:
la aceleración no altera el formato de los \textit{checkpoints} ni
la reanudación de entrenamientos.

\paragraph{Validación durante el entrenamiento.}
Como se observó en el experimento de mariposas
(\Cref{sec:mariposas-aprendizajes}), la pérdida de entrenamiento es
una mala señal de convergencia en difusión: el MSE sobre
$\mathbf{v}$, promediado sobre \textit{timesteps} aleatorios, se
satura en las primeras épocas mientras la calidad de imagen sigue
mejorando. Por ello, en cada \textit{checkpoint} se ejecuta una
validación \emph{de tarea final}: se reconstruye un subconjunto fijo
de 16 cortes de validación —equiespaciados sobre todo el split para
cubrir distintos niveles anatómicos— con el mismo procedimiento de
inferencia que se usará en test, usando los pesos EMA, y se
registran PSNR/SSIM/NRMSE contra el objetivo de dosis completa. La
convergencia se declara cuando estas métricas se estancan. El
generador aleatorio de la validación se aísla del de entrenamiento
para no perturbar la reproducibilidad del bucle principal.

Los \textit{checkpoints} se guardan cada 10 épocas en directorios
autocontenidos (pesos EMA del U-Net en formato \texttt{diffusers},
configuración del \textit{scheduler} y estado completo de
entrenamiento), lo que permite tanto la inferencia directa como la
reanudación exacta del entrenamiento tras una interrupción
—condición imprescindible cuando se alquila cómputo por horas—.

\subsection{Entorno de ejecución en la nube}
\label{sec:pet-entorno-cloud}

\subsubsection{Motivación: del portátil a la GPU dedicada}

El plan inicial contemplaba ejecutar el entrenamiento en local,
sobre el mismo MacBook Pro M4 Pro (24~GB de memoria unificada,
PyTorch con \textit{backend} MPS) del experimento de mariposas. Las
estimaciones lo descartaron: en MPS, forzado a \texttt{float32} y
con \textit{batch} efectivo 16, el coste estimado era de
aproximadamente \textbf{75~horas por \textit{pipeline}} de difusión,
además de cerca de 2~horas por evaluación completa del conjunto de
test. Más allá del tiempo bruto, 150~horas de entrenamiento
sostenido acoplan el ordenador personal al proceso durante más de
una semana, con cada interrupción (cierre de tapa, suspensión,
reinicio) convertida en riesgo operativo. Se decidió migrar la fase
PET a \textbf{RunPod}, un proveedor de GPU bajo demanda con tarifa
proporcional al uso, y construir alrededor un entorno de ejecución
reproducible cuyo diseño se describe a continuación.

\subsubsection{Imagen Docker reproducible}

Todo el entorno de software se encapsula en una imagen Docker
construida sobre \texttt{python:3.13-slim}. Dos decisiones de diseño
merecen comentario:

\begin{itemize}[leftmargin=*]
    \item \textbf{CUDA sin imagen base de NVIDIA.} Los
    \textit{wheels} de PyTorch del índice \texttt{cu124} empaquetan
    las bibliotecas CUDA necesarias, de modo que basta una imagen
    base ligera de Python: el único requisito del host es el driver
    NVIDIA, que RunPod provee. PyTorch se instala antes que el resto
    de dependencias para que \texttt{pip} no resuelva por error la
    variante de CPU, y la capa de dependencias se cachea separada
    del código para acelerar las reconstrucciones de la imagen.
    \item \textbf{Compilador de C en la imagen.} TorchInductor
    (\texttt{torch.compile}, \Cref{subsec:pet-config}) genera y
    compila kernels en tiempo de ejecución y necesita \texttt{gcc};
    se incluye \texttt{build-essential} en la imagen, sin lo cual la
    compilación falla al arrancar el entrenamiento.
\end{itemize}

La imagen \emph{no} hornea el código del proyecto: un
\textit{entrypoint} idempotente lo clona desde GitHub al arrancar el
contenedor, de modo que cada arranque parte del repositorio
actualizado sin reconstruir la imagen. El mismo \textit{entrypoint}
arranca un servidor SSH —inyectando la clave pública que RunPod
expone por variable de entorno— antes de cualquier paso costoso,
para poder conectarse mientras se prepara el entorno. Un fichero
\texttt{docker-compose} permite validar la imagen y el
\textit{entrypoint} en local (emulando \texttt{linux/amd64}) antes
de desplegar.

\subsubsection{Distribución del dataset preprocesado vía Hugging Face}

Subir 60~GB de caché preprocesada a cada máquina efímera habría sido
el cuello de botella del flujo. En su lugar, la caché completa
(\Cref{sec:pet-preprocesado}) se empaqueta una única vez en
\textit{shards} de \texttt{tar} y se publica —junto al fichero de
particiones \texttt{splits.json}— en un repositorio \emph{privado}
de Hugging Face Datasets. Al arrancar, el \textit{entrypoint}
descarga los \textit{shards} con \texttt{hf\_transfer} (descarga
paralela de alto ancho de banda), reconstruye la caché, y
\textbf{verifica su integridad} comparando el número de ficheros
\texttt{.pt} reconstruidos (477.690) con el esperado antes de borrar
los \textit{shards}; si la cuenta no coincide, conserva la descarga
para depuración y deja el contenedor vivo. Una marca en disco hace
el proceso idempotente: los reinicios posteriores del pod detectan
el dataset ya reconstruido y lo omiten.

\subsubsection{Elección de almacenamiento y de máquina}

RunPod separa el cómputo del almacenamiento:

\begin{description}[leftmargin=*]
    \item[\textit{Network volume} (almacenamiento persistente).] Se
    asocia al pod un volumen de red que sobrevive al apagado de la
    instancia. Sobre él se materializan el repositorio, la caché de
    \textit{slices} ($\sim 60$~GB), los \textit{checkpoints} y los
    \textit{logs} de TensorBoard. De esta manera, detener la GPU
    para no pagar por horas ociosas no implica perder el estado del
    trabajo ni volver a descargar el dataset; el volumen se factura
    por GB-mes, una fracción pequeña del coste de GPU
    (\Cref{sec:pet-costes}).
    \item[Instancia GPU.] Se selecciona una \textbf{NVIDIA RTX PRO
    4500 (arquitectura Blackwell, 32~GB de VRAM)}. La elección
    responde a tres criterios: (i)~\emph{memoria}: 32~GB permiten el
    \textit{batch} de 32 cortes a $256^2$ con el U-Net de 114~M de
    parámetros más su copia EMA sin acumulación de gradiente;
    (ii)~\emph{generación}: los \textit{tensor cores} Blackwell
    rinden el \texttt{bf16} y el TF32 de los que depende la
    configuración de la \Cref{subsec:pet-config}; y
    (iii)~\emph{precio}: las alternativas de consumo con menos VRAM
    (RTX 4090, 24~GB) habrían obligado a acumulación de gradiente.
    Esta tarjeta fija el \emph{suelo} de hardware del proyecto —los
    32~GB de VRAM para los que se calibró la configuración de la
    \Cref{subsec:pet-config}— y la máquina de referencia sobre la
    que se desarrolla y depura el código; recupera la precisión mixta
    y el \textit{batch} completo, y desacopla por completo el
    entrenamiento del ordenador personal. Ahora bien, RunPod es un
    mercado de capacidad bajo demanda y la disponibilidad real de
    cada modelo de tarjeta fluctúa, de modo que los experimentos
    acabaron ejecutándose sobre un \textit{pool} rotatorio de GPUs
    —incluida la GPU de centro de datos A100 cuando estaba
    disponible a buen precio, para acortar el tiempo de respuesta y
    paralelizar entrenamientos—, como se detalla a continuación.
\end{description}

\subsubsection{Operativa de ejecución}

El flujo de trabajo sobre el pod descansa en herramientas estándar
de administración de sistemas Linux remotos:

\begin{description}[leftmargin=*]
    \item[SSH.] La conexión al pod se realiza por SSH sobre la
    dirección y puerto que RunPod expone tras lanzar la instancia,
    autenticada con la clave pública inyectada por el
    \textit{entrypoint}.
    \item[\texttt{tmux}.] Cada experimento se lanza dentro de una
    sesión \texttt{tmux} en el pod. \texttt{tmux} desacopla el ciclo
    de vida del proceso del de la conexión SSH: cerrar la terminal,
    perder la red o reiniciar el portátil no interrumpen el
    entrenamiento, y reconectarse permite recuperar la sesión
    exactamente en el estado en que se dejó. Combinado con la
    reanudación desde \textit{checkpoint}
    (\Cref{subsec:pet-config}), el entrenamiento sobrevive tanto a
    desconexiones como a reinicios del propio pod.
    \item[Editor remoto.] La edición y depuración puntual del código
    se realiza con Cursor (un \textit{fork} de Visual Studio Code) y
    la extensión \textit{Remote--SSH}: el editor corre en local pero
    opera sobre el sistema de archivos del pod por el mismo canal
    SSH, sin sincronizar copias del código.
\end{description}

El código del proyecto está organizado para que esta operativa sea
de un solo comando: un \emph{dispatcher} único
(\texttt{python -m src.main \{preprocess|train|reconstruct|evaluate|preview\}})
enruta cada subcomando al módulo correspondiente, con un flag
\texttt{--pipeline} que selecciona entre los dos \textit{pipelines}
de difusión y las dos líneas base, y un preset \texttt{--smoke}
(resolución $128^2$, 5 épocas, subconjunto de pacientes) que
permite verificar el flujo completo de extremo a extremo en
aproximadamente una hora antes de comprometer un entrenamiento
largo. La infraestructura compartida (motor de entrenamiento,
validación, E/S de volúmenes, métricas) vive en módulos comunes, de
modo que los cuatro experimentos difieren únicamente en la
construcción de la entrada del modelo y en su procedimiento de
inferencia.

\subsubsection{\textit{Pool} de GPUs y ejecución en paralelo}
\label{subsec:pet-pool-gpu}

La combinación de volumen persistente y reanudación desde
\textit{checkpoint} (\Cref{subsec:pet-config}) convierte la GPU en
un recurso intercambiable: detener un pod, lanzar otro con una
tarjeta distinta y reanudar el entrenamiento es transparente para el
estado del experimento. Esto permitió aprovechar el mercado bajo
demanda de RunPod y asignar a cada experimento la tarjeta más
conveniente entre las disponibles en cada momento. El conjunto de
GPUs empleado se recoge en la \Cref{tab:pet-gpus}.

\begin{table}[h!]
    \centering
    \small
    \begin{tabular}{lccc}
        \toprule
        GPU & VRAM & vCPU & Tarifa (USD/h) \\
        \midrule
        NVIDIA A100 SXM      & 80~GB & 16 & $1{,}49$ \\
        NVIDIA RTX PRO 4500  & 32~GB & 12 & $0{,}74$ \\
        NVIDIA RTX A6000     & 48~GB &  9 & $0{,}49$ \\
        NVIDIA A40           & 48~GB &  9 & $0{,}44$ \\
        \bottomrule
    \end{tabular}
    \caption{\textit{Pool} de instancias GPU de RunPod empleadas, con
    su tarifa horaria. La configuración de entrenamiento se calibró
    para el suelo de 32~GB (RTX PRO 4500), de modo que corre sin
    cambios en cualquiera de ellas; las tarjetas con más VRAM se
    reservaron para acelerar y paralelizar los entrenamientos más
    costosos.}
    \label{tab:pet-gpus}
\end{table}

Como la configuración exige solo 32~GB, las tarjetas con más memoria
no aportan \textit{batch} adicional, pero sí más rendimiento de
cómputo —y, sobre todo, la posibilidad de \textbf{entrenar dos
\textit{pipelines} simultáneamente} en dos instancias independientes,
reduciendo a la mitad el tiempo de pared para completar los dos
modelos de difusión—. La \Cref{fig:pet-gpu-util} muestra el panel de
monitorización de RunPod durante una de esas ejecuciones en paralelo,
con los dos entrenamientos de difusión corriendo a la vez sobre
sendas instancias A100 cercanas a la saturación de cómputo.

\begin{figure}[h!]
    \centering
    \includegraphics[width=\linewidth]{img/GPU-utilization.png}
    \caption{Panel de monitorización de RunPod durante el
    entrenamiento simultáneo de los dos \textit{pipelines} de
    difusión sobre dos instancias A100 independientes. Ambos pods
    operan con la GPU prácticamente saturada (utilización de cómputo
    en torno al 93--100\,\%) y un uso de VRAM holgado —la
    configuración se calibró para 32~GB, muy por debajo de los 80~GB
    de la A100—, lo que confirma que el cuello de botella es el
    cómputo y no la memoria.}
    \label{fig:pet-gpu-util}
\end{figure}

% ---------------------------------------------------------------------
\section{Arquitectura y enfoque}
\label{sec:pet-arquitectura}

\subsection{Red de \textit{denoising}}

La arquitectura del U-Net es idéntica a la empleada en el modelo
de mariposas (\Cref{sec:mariposas-arquitectura}), modificando
exclusivamente los canales de entrada y salida:

\begin{itemize}[leftmargin=*]
    \item \texttt{block\_out\_channels = (128, 128, 256, 256, 512,
    512)}, lo cual produce un modelo de aproximadamente 114
    millones de parámetros (idéntico recuento al modelo de
    mariposas, dado que el primer canal convolucional contribuye
    una fracción despreciable al total).
    \item \texttt{layers\_per\_block = 2}.
    \item Bloque de atención multi-cabeza en el quinto nivel del
    camino descendente y ascendente (resolución espacial
    $16 \times 16$ tras cinco \textit{downsamplings} del
    \textit{input} de $256 \times 256$).
    \item El \textit{timestep} $t$ se inyecta mediante
    \textit{embeddings} sinusoidales en cada bloque ResNet.
\end{itemize}

Los canales de entrada y salida dependen del \textit{pipeline}:
\texttt{in\_channels = 2} en el supervisado —canal de difusión
$\mathbf{x}_t$ y canal de condicionamiento $\mathbf{y}$— y
\texttt{in\_channels = 1} en el no condicional; en ambos casos
\texttt{out\_channels = 1}.

\subsection{Pipeline~A — supervisado condicional}
\label{sec:pet-pipeline-a}

El Pipeline~A instancia el condicionamiento supervisado de la
\Cref{subsec:difusion-condicional} mediante \textbf{concatenación
por canales}, el mecanismo dominante en la literatura supervisada
de PET de dosis baja con DDPMs. La entrada del U-Net es el tensor
de forma $2 \times 256 \times 256$ obtenido concatenando
$\mathbf{x}_t$ —el \textit{slice} de dosis completa con ruido
inyectado en el \textit{timestep}~$t$ según la
\Cref{eq:forward-reparam}— y $\mathbf{y}$ —el \textit{slice}
pareado de dosis baja, normalizado a la misma escala que
$\mathbf{x}_0$—. La salida es la predicción $\hat{\mathbf{v}}_t$ de
tamaño $1 \times 256 \times 256$, entrenada con la
\Cref{eq:loss-v}.

En inferencia, para cada \textit{slice} axial del volumen se parte
de $\mathbf{x}_T \sim \mathcal{N}(\mathbf{0}, \mathbf{I})$
concatenado con el correspondiente $\mathbf{y}$, y se aplica el
muestreador DDIM (\Cref{eq:ddim-step}) con 50 pasos de
\textit{denoising} y $\eta = 0$, lo cual hace el proceso
determinista para una semilla dada.

\subsection{Pipeline~B — prior no condicional con guiado de medición
(DPS)}
\label{sec:pet-pipeline-b}

El entrenamiento del Pipeline~B es idéntico al del Pipeline~A
pero \textbf{sin el canal de condicionamiento}: el U-Net aprende
$p(\mathbf{x})$ sobre los \textit{slices} de dosis completa
exclusivamente, y la observación $\mathbf{y}$ no interviene en el
entrenamiento. En inferencia se introduce la consistencia con la
medición mediante \textit{Diffusion Posterior Sampling}
(\Cref{subsec:difusion-condicional}): en cada paso DDIM, la red
predice $\hat{\mathbf{v}}_t$, se reconstruye la estimación de
imagen limpia $\hat{\mathbf{x}}_0$ con la \Cref{eq:x0-from-v}, se
aplica el paso de gradiente de consistencia de la
\Cref{eq:dps-update} con operador directo identidad ($A =
\mathbf{I}$), y se ejecuta el paso DDIM estándar sobre el estado ya
corregido.

La elección del operador identidad y de una varianza de observación
uniforme es la aproximación más simple y defendible en este
contexto: el PET de dosis baja carece de una formulación limpia
$\mathbf{y} = A\mathbf{x}$ —la reducción de dosis combina
re-muestreo de Poisson, reducción del tiempo de escaneo y la
posterior etapa de reconstrucción algorítmica del escáner, ninguna
de las cuales se modela explícitamente—. La pérdida cuadrática en
el espacio normalizado por $\operatorname{arcsinh}$ equivale a una
verosimilitud gaussiana isotrópica sobre la representación
estabilizada de varianza. La sustitución de la varianza uniforme
por una $\sigma(\mathbf{y})$ derivada de la estadística de Poisson
es una de las extensiones identificadas como líneas de desarrollo
futuro (\Cref{sec:lineas-futuras}). La escala de guiado $\omega$ es
el único hiperparámetro adicional del procedimiento, y se ajusta
sobre el conjunto de validación.

\subsection{Inferencia volumétrica y exportación}
\label{subsec:pet-inferencia-volumetrica}

En ambos \textit{pipelines}, la reconstrucción de un volumen lee el
NIfTI de dosis baja \emph{crudo} y deriva en tiempo de ejecución
—con la misma función que usa el preprocesamiento— la
\textit{bounding box}, la escala $M$ y los cortes a procesar, sin
consultar ninguna estadística precalculada. Esta decisión tiene una
consecuencia operativa importante: el sistema funciona sobre un
volumen \emph{no visto} que nunca pasó por el preprocesado, igual
que ocurriría en un despliegue real. Tras muestrear todos los
\textit{slices}, éstos se reensamblan en un volumen tridimensional,
se invierte la normalización de la \Cref{eq:asinh-norm}, se deshace
el reescalado y el recorte, y el resultado se escribe en formato
NIfTI con la matriz afín del volumen de entrada.

% ---------------------------------------------------------------------
\section{Métricas de evaluación}
\label{sec:pet-metricas}

Sobre el conjunto de test (38 volúmenes) se reportan tres familias
de métricas.

\paragraph{Métricas de calidad de imagen.}
PSNR, SSIM y NRMSE se calculan por \textit{slice} y se agregan
por volumen mediante la media a lo largo de los \textit{slices}.
Se reportan dos variantes: (i) sobre el \textit{slice} completo
y (ii) restringidas a la máscara de \textit{foreground} del
\textit{slice}. La variante restringida es la clínicamente
relevante porque ignora el fondo vacío; la variante completa es
la convención estándar y permite la comparación con la literatura.

Estas tres métricas se calculan sobre el \textbf{espacio
normalizado} de las imágenes (la representación
$\operatorname{arcsinh}$ en $[-1, +1]$), y no sobre el espacio de
cuentas. La razón es que el PET tiene un rango dinámico altísimo:
el PSNR colapsa hacia el vóxel más brillante, por lo que evaluarlo
en cuentas equivaldría a medir casi exclusivamente cómo de bien se
reproduce el punto más caliente, siendo numéricamente insensible
a los errores del resto de la imagen; análogamente, los valores
atípicos inflarían el SSIM. Además, promediar a lo largo de
pacientes requiere una escala común para que las comparaciones
intra-volumen sean justas. La reformulación clave es que
$\operatorname{arcsinh}$ es monótona e invertible y se aplica por
igual a la reconstrucción y al \textit{ground truth}: no oculta el
error, sino que lo repondera para que cuenten los errores
relativos en todo el rango dinámico en lugar de los errores
absolutos en el pico. Así, PSNR/SSIM en espacio normalizado son
una comparación fiel del resultado final, bajo una ponderación de
intensidad en la que la métrica no está degenerada.

\paragraph{Preservación de intensidad.}
Como métrica \textit{proxy} de la preservación de SUV —no se
dispone de metadatos de dosis inyectada ni peso del paciente para
computar SUV reales—, se reporta el error relativo entre la
\textbf{media} y el \textbf{máximo} de intensidad de
\textit{foreground} por volumen, calculado en \textbf{espacio
de cuentas original} (tras invertir la \Cref{eq:asinh-norm} con la
$M$ y $k$ derivadas de la dosis baja del propio volumen).

\paragraph{Evaluación cualitativa.}
Para $N = 5$ volúmenes representativos del conjunto de test se
generan figuras de cuatro paneles (dosis baja, dosis completa,
reconstrucción, residuo absoluto) que permiten un análisis visual
del comportamiento del modelo.

Todas las comparaciones de la \Cref{sec:pet-resultados} se realizan
sobre exactamente los mismos volúmenes y los mismos
\textit{slices}, con idéntica normalización per-volumen derivada de
la dosis baja, lo cual garantiza que las métricas son directamente
comparables entre métodos.

% ---------------------------------------------------------------------
\section{Líneas base de referencia}
\label{sec:pet-baselines}

Comparar los dos \textit{pipelines} de difusión solo entre sí
dejaría una pregunta esencial sin responder: ¿cuánto de la mejora se
debe al paradigma de difusión y cuánto se obtendría con métodos más
simples? Para delimitarlo se evalúan tres referencias, ordenadas de
menor a mayor capacidad.

\paragraph{Dosis baja sin procesar.}
La referencia inferior es la propia observación $\mathbf{y}$,
evaluada directamente contra la dosis completa con las métricas de
la \Cref{sec:pet-metricas}. Es el suelo que cualquier método debe
superar, y cuantifica la severidad de la degradación de partida.

\paragraph{Línea base A — U-Net de regresión (ablación del
paradigma).}
La \emph{misma} U-Net del \textit{pipeline} supervisado —idénticos
bloques, canales y número de parámetros— entrenada como regresor
directo: entrada de un solo canal (el corte de dosis baja), salida
de un canal (la estimación de dosis completa), objetivo MSE en el
espacio normalizado $[-1, 1]$, sin proceso de difusión ni
\textit{scheduler} de ruido. La inferencia es una única pasada hacia
delante. Mismos datos, misma normalización y mismo presupuesto de
entrenamiento que el Pipeline~A: la comparación aísla exactamente la
aportación del \emph{paradigma generativo}, manteniendo constante
todo lo demás.

\paragraph{Línea base B — RED-CNN (ablación de la arquitectura).}
RED-CNN~\cite{redcnn2017} es la CNN clásica de referencia para
\textit{denoising} de TAC y PET de dosis baja: un
codificador-decodificador residual totalmente convolucional de 5
capas de convolución y 5 de deconvolución, 96 filtros de
$5 \times 5$, paso 1 y sin \textit{pooling} —opera a resolución
completa—, con conexiones residuales simétricas. Son
$\sim 1{,}8$ millones de parámetros, unas 60 veces menos que la
U-Net. Se entrena de forma idéntica a la línea base~A (mismos pares,
mismo objetivo MSE en espacio normalizado, mismo presupuesto), de
modo que la comparación con ella aísla la \emph{familia
arquitectónica} en lugar del paradigma. La única adaptación respecto
al original es la salida lineal (se elimina la ReLU final), ya que
el espacio normalizado es con signo y el fondo vive en $-1$.

% ---------------------------------------------------------------------
\section{Comparación de resultados}
\label{sec:pet-resultados}

La \Cref{tab:pet-comparativa} recoge las métricas de los cuatro
modelos y de la referencia de dosis baja sobre el conjunto de test
completo (38 volúmenes, 24.472 \textit{slices}), evaluadas conforme
a la \Cref{sec:pet-metricas}.

\begin{table}[h!]
    \centering
    \small
    \begin{tabular}{lccccc}
        \toprule
        Métrica & Dosis baja & RED-CNN & U-Net reg. &
        \begin{tabular}{@{}c@{}}Difusión\\supervisada\end{tabular} &
        \begin{tabular}{@{}c@{}}Difusión\\incond.\ (DPS)\end{tabular} \\
        \midrule
        PSNR \textit{fg} (dB) $\uparrow$ & $24{,}56$ & \todo{} & \todo{} & $30{,}20$ & \todo{} \\
        PSNR \textit{whole} (dB) $\uparrow$ & $28{,}31$ & \todo{} & \todo{} & $35{,}32$ & \todo{} \\
        SSIM \textit{fg} $\uparrow$ & $0{,}9201$ & \todo{} & \todo{} & $0{,}9652$ & \todo{} \\
        SSIM \textit{whole} $\uparrow$ & $0{,}7308$ & \todo{} & \todo{} & $0{,}9251$ & \todo{} \\
        NRMSE \textit{fg} $\downarrow$ & $0{,}0701$ & \todo{} & \todo{} & $0{,}0364$ & \todo{} \\
        NRMSE \textit{whole} $\downarrow$ & $0{,}0434$ & \todo{} & \todo{} & $0{,}0194$ & \todo{} \\
        Error medio intensidad (\%) & $+1{,}77$ & \todo{} & \todo{} & $-4{,}75$ & \todo{} \\
        Error máx.\ intensidad (\%) & $+50{,}10$ & \todo{} & \todo{} & $-18{,}47$ & \todo{} \\
        \bottomrule
    \end{tabular}
    \caption{Comparativa sobre el conjunto de test (38 volúmenes,
    24.472 \textit{slices}). PSNR/SSIM/NRMSE en espacio normalizado
    (\textit{fg} = restringido a \textit{foreground};
    \textit{whole} = \textit{slice} completo); errores de intensidad
    en espacio de cuentas. La columna de difusión supervisada
    corresponde al \textit{checkpoint} de la época~100.}
    \label{tab:pet-comparativa}
\end{table}

\todo{Completar las columnas de RED-CNN, U-Net de regresión y
difusión incondicional (DPS) cuando termine la recopilación de
resultados de los últimos experimentos, e insertar las figuras
representativas de cuatro paneles.}

Los resultados ya disponibles permiten dos lecturas.

\paragraph{Difusión supervisada frente a la observación.}
El Pipeline~A mejora a la dosis baja sin procesar en todas las
métricas estructurales: $+7{,}0$~dB de PSNR sobre el \textit{slice}
completo ($28{,}31 \to 35{,}32$), $+5{,}6$~dB en \textit{foreground},
y un salto del SSIM \textit{whole} de $0{,}73$ a $0{,}93$. El NRMSE
se reduce aproximadamente a la mitad en ambas variantes. En
preservación de intensidad la lectura es más matizada: la dosis baja
preserva casi exactamente la \emph{media} de intensidad por
construcción (error $+1{,}8$\,\%) pero distorsiona gravemente el
\emph{máximo} ($+50$\,\%, el ruido de Poisson infla los picos),
mientras que la reconstrucción supervisada estabiliza el máximo
($-18{,}5$\,\%) a costa de una ligera subestimación sistemática de
la señal media ($-4{,}75$\,\%), coherente con la saturación suave
esperable en las regiones más calientes
(\Cref{sec:pet-preprocesado}).

\paragraph{Progresión del entrenamiento.}
Entre los \textit{checkpoints} de las épocas 50 y 100 del
Pipeline~A, las métricas de calidad de imagen mejoran de forma
consistente (PSNR \textit{whole} $34{,}52 \to 35{,}32$~dB; SSIM
\textit{whole} $0{,}9207 \to 0{,}9251$; NRMSE \textit{whole}
$0{,}0213 \to 0{,}0194$), confirmando la observación del
\Cref{cap:mariposas} de que la calidad sigue creciendo mucho después
de que la pérdida de entrenamiento se estabilice. El error de
intensidad media, en cambio, empeora ligeramente
($-4{,}18 \to -4{,}75$\,\%), lo cual sugiere que el óptimo de las
métricas estructurales y el de la fidelidad cuantitativa no
coinciden necesariamente en la misma época.

% ---------------------------------------------------------------------
\section{Costes}
\label{sec:pet-costes}

Uno de los objetivos declarados del trabajo
(\Cref{sec:objetivos}) es documentar el coste real de los
experimentos, una información que la literatura rara vez hace
explícita y que condiciona directamente qué se puede y qué no se
puede explorar en un proyecto con presupuesto acotado.

El modelo de facturación de RunPod tiene dos componentes:

\begin{itemize}[leftmargin=*]
    \item \textbf{Cómputo}: cada instancia GPU del \textit{pool}
    (\Cref{tab:pet-gpus}) se factura por horas de pod activo a su
    tarifa correspondiente. Detener el pod detiene el gasto de
    cómputo, y la combinación de volumen persistente + reanudación
    desde \textit{checkpoint} (\Cref{sec:pet-entorno-cloud}) permite
    hacerlo sin perder estado, de modo que solo se paga por horas de
    GPU realmente útiles.
    \item \textbf{Almacenamiento}: el \textit{network volume} se
    factura por GB-mes con independencia de que el pod esté activo.
    Con $\sim$100~GB (caché de 60~GB más \textit{checkpoints} y
    \textit{logs}) su coste mensual es una fracción pequeña del
    coste horario de la GPU.
\end{itemize}

El alojamiento del dataset preprocesado en Hugging Face Datasets es
gratuito, y la transferencia de datos de entrada está incluida en la
tarifa del pod. La \Cref{tab:pet-costes} desglosa el coste de los
cuatro experimentos de entrenamiento, asignando a cada uno la GPU del
\textit{pool} (\Cref{tab:pet-gpus}) sobre la que se ejecutó.

\begin{table}[h!]
    \centering
    \small
    \begin{tabular}{llccc}
        \toprule
        Experimento & GPU & Horas & Tarifa (USD/h) & Coste (USD) \\
        \midrule
        Difusión supervisada (Pipeline A)   & A100 SXM     & $38$ & $1{,}49$ & $56{,}62$ \\
        Difusión incondicional (Pipeline B) & A100 SXM     & $40$ & $1{,}49$ & $59{,}60$ \\
        U-Net de regresión (línea base A)   & RTX A6000    & $34$ & $0{,}49$ & $16{,}66$ \\
        RED-CNN (línea base B)              & A40          & $18$ & $0{,}44$ & $7{,}92$ \\
        \midrule
        Evaluaciones e inferencia           & RTX PRO 4500 & $22$ & $0{,}74$ & $16{,}28$ \\
        \textit{Network volume} ($\sim$100~GB-mes) & --- & \multicolumn{2}{c}{---} & $11{,}00$ \\
        \midrule
        \textbf{Total} & & $\mathbf{152}$ & & $\mathbf{168{,}08}$ \\
        \bottomrule
    \end{tabular}
    \caption{Coste de los experimentos sobre RunPod, con la GPU del
    \textit{pool} (\Cref{tab:pet-gpus}) empleada en cada uno. Cada
    fila de entrenamiento corresponde a 100 épocas sobre los
    $\sim 1{,}9 \times 10^{5}$ pares de cortes de entrenamiento. Los
    dos \textit{pipelines} de difusión se ejecutaron sobre A100 —en
    parte de forma simultánea (\Cref{fig:pet-gpu-util})— para
    acortar el tiempo de pared; las dos líneas base discriminativas,
    más ligeras, se relegaron a las tarjetas más económicas del
    \textit{pool}.}
    \label{tab:pet-costes}
\end{table}

El coste agregado de cómputo documentado asciende a
\textbf{$\approx 168$~USD} sobre 152~horas de GPU, repartidas entre
las cuatro tarjetas del \textit{pool}. Las dos líneas discriminativas
ilustran bien el desacople entre coste y arquitectura: RED-CNN, con
$\sim 1{,}8$~millones de parámetros, completa su presupuesto de 100
épocas en menos de la mitad de horas que la U-Net de regresión pese a
operar a resolución completa, y sobre la tarjeta más barata del
conjunto. El saldo total aprovisionado en RunPod a lo largo del
proyecto fue algo mayor que esta cifra —se recargó en incrementos
pequeños según se necesitaba— y absorbió además las pruebas
\texttt{-{}-smoke}, la depuración interactiva, algún relanzamiento
tras corregir errores y el inevitable tiempo de pod ocioso antes de
detener cada instancia, además de un margen que quedó sin consumir.

La comparación con la estimación inicial es la lección económica del
capítulo. Ejecutar los dos \textit{pipelines} de difusión en el
portátil se había estimado en torno a 150~horas de cómputo acoplado
al ordenador personal (\Cref{sec:pet-entorno-cloud}), inviable en la
práctica por el riesgo operativo de mantener la máquina dedicada
durante más de una semana. Migrar a GPU dedicada no solo eliminó ese
acoplamiento, sino que redujo el tiempo de pared a unas pocas
decenas de horas por modelo —y a la mitad cuando se entrenaron dos a
la vez— por un desembolso monetario modesto, plenamente dentro del
presupuesto de un Trabajo de Fin de Máster. El cómputo dejó de ser,
así, el factor limitante del alcance experimental.

Más allá del entrenamiento, conviene registrar el coste de
\emph{inferencia}: reconstruir un volumen completo con DDIM-50 en
local (MacBook Pro M4 Pro, \textit{batch} de inferencia 8) requiere
aproximadamente una hora por volumen, lo que sitúa la inferencia de
difusión lejos del tiempo real y motiva la discusión sobre el
despliegue clínico de la \Cref{sec:limitaciones}. Las líneas base
discriminativas, con una única pasada hacia delante por corte, son
órdenes de magnitud más rápidas en inferencia; este contraste forma
parte de la comparación de paradigmas y se retoma en las
conclusiones.
